% Pursuing Stacks (cotes 134-2 to 134-8): archive cote 134-4, site page 9 (PDF page 10), typescript page 97.
% typescript. Draft compared: ps3.tex:419-491.
% AI-transcribed by Claude Opus 5.5 (claude-opus-5-5, adaptive thinking, effort high), Message Batches API, batch ps1 (msgbatch_01Rv5hCaNTtUt4JHx1yqoHPL), 2026-11-12; not yet checked by eye.
\scanpage{134-4/9}
\mspage{97}

still, with (Cat) replaced by such an M.

In the situation here, the work will consist in devising handy condi-%
\inserted[on i: $A\rightarrow$ (Cat)]{and A,}
tions which will ensure that all three maps in (1) are weak equivalences,
for any choice of C. This will imply that the corresponding maps in (Hot)
are isomorphisms, hence a canonical isomorphism between the images in
(Hot) of $A_{/C}$ and C , which will imply that a) the functor $i_A i^*$ from
(Cat) to (Cat) carries weak equivalences into weak equivalences, and
hence induces a functor
\[ (Hot) \longrightarrow (Hot) , \]
and b) that this functor is \deleted{equivalent}\inserted{isomorphic} to the identity functor. If
moreover A is a weak test category, and therefore the functor
\[ W_A^{-1}A^{\wedge} \longrightarrow (Hot) \]
induced by $i_A$ is an equivalence, it will follow that \xout{induced a}\inserted[the]{} functor
\[ (Hot) \longrightarrow W_A^{-1}A^{\wedge} \]
induced by $i^*$ is equally an equivalence, namely that $i^*$ is indeed a
test-functor.

For the map
\[ (2) \qquad AxC \rightarrow C \]
to be a weak equivalence for any C in (Cat), it is necessary and suffi-%
cient that A be aspheric (cf. par.40, page 69), a familiar condition on
A indeed ! For handling the map
\[ (3) \qquad AxC \rightarrow A_{//C} \qquad \text{associated to} \quad C \rightarrow \underline{Hom}(i(a),C) \ , \]
we'll use the following easy result (which I'll admit for the time being):
\uline{Proposition} Let F and G be two \xout{fib} categories over a category A, and
u: $F \rightarrow G$ a functor compatible with projections. We assume
a) For any a in A, the induced map on the fibers
\[ u_a : F_a \rightarrow G_a \]
is a weak equivalence.
b) Either F and G are both cofibering over A and u is cocartesian,
or F and G are fibering and i is cartesian.
Then u is a weak equivalence.

This shows that a sufficient condition for (3) to be a weak equivalen-%
ce, is that the functor\xout{s} $C \rightarrow \underline{Hom}(i(a),C)$ be a weak equivalence, for
any a in A and C in (Cat). We'll see in the next section\hand{*} that this amount\xout{\unsure{s}}
to demanding that the \xout{categories} objects i(a) in (Cat) should be ``contrac%
% NOTE: the right edge of the scan cuts this line after "contrac"; a hyphen may be lost
tible'', in the most concrete sense of this \xout{exp} expression, which is ac-%
tually stronger than just asphericity. \hand{(}Earlier in these notes I was a \margin{*9.4. Actually this is done only a lot later, on page 121 and (for the converse) on page 143.}
little floppy with the termin\inserted{o}logy, by using a few times the word ``con-%

\margin{9}
